\usepackage{xspace}
\newtheorem{definition}{Definition}

\newcommand{\hua}[1]{\textcolor{purple}{\bf [#1 -Hua]}}
\newcommand{\chuke}[1]{\textcolor{green}{\bf [#1 -Chuke]}}
\newcommand{\kaifang}[1]{\textcolor{orange}{\bf [#1 -Kaifang]}}

\newcommand{\system}
{{\textsc{HalluciActionLens}}\xspace}

% Severity highlight: the tinted themes in Table 2 and the "Red" swatch in its
% caption.  Defined once, here, for two reasons.
%
% \fboxsep has to be set and it has to be small.  At the LaTeX default of 3pt
% the box stands 3pt above and below a normal line, and in running text that is
% enough to break the leading: the line no longer fits under \baselineskip, TeX
% falls back to \lineskip, and the highlighted line drops ~2pt away from its
% neighbours while every other line keeps its spacing.  1.4pt clears a 10pt/11pt
% line with room to spare.  (This is what the caption of Table 2 used to do --
% \fboxsep was set to 1.4pt inside tables/rq3_failure_modes, which is \input
% *after* the caption, so the caption itself ran at the 3pt default.)
%
% And it keeps the two uses identical.  The caption says "Red marks a theme",
% so its swatch has to be the same size and colour as the marks it describes;
% with the padding set in one place and the other left at the default they were
% visibly different boxes.
%
% The braces scope \fboxsep, which the old table-local \setlength did not -- it
% leaked 1.4pt into every \fbox and \colorbox in the rest of the document.
\newcommand{\hlred}[1]{{\setlength{\fboxsep}{1.4pt}\colorbox{red!16}{#1}}}
